\begin{ouverture}[calcul]{1}{Entiers et décimaux}{Partie 1 · Nombres et calculs}
{Milliards · Unités de numération · Écritures d’un nombre}
\objectif{Lire, écrire et décomposer de grands nombres, jusqu’aux milliards.}
\objectif{Connaître les liens entre unités, dixièmes, centièmes et millièmes.}
\objectif{Reconnaître un nombre décimal et passer d’une écriture à une autre.}
\prerequis{Numération du CM2 : entiers jusqu’aux centaines de millions, décimaux jusqu’aux millièmes.}
\end{ouverture}
\pagestyle{otMain}\essentiel\fontsize{11}{15}\selectfont
\phantomsection\label{cours-1-1}
\essSection{1}{Les grands nombres et les milliards}
\begin{definition}Les chiffres d’un entier se regroupent par classes de trois, de droite à gauche : unités, mille, millions, milliards. Un milliard, c’est mille millions : $1\,000\,000\,000$. Le nombre $7\,342\,015\,608$ se lit « sept-milliards-trois-cent-quarante-deux-millions-quinze-mille-six-cent-huit ».\end{definition}
\begin{center}\renewcommand{\arraystretch}{1.35}\begin{tabular}{c|cccc}\toprule
Classe & milliards & millions & mille & unités\\\midrule
Rangs & c d u & c d u & c d u & c d u\\
Chiffres & 7 & 3 4 2 & 0 1 5 & 6 0 8\\
\bottomrule\end{tabular}\end{center}
\begin{savoirfaire}[Chiffre ou nombre de]Dans $7\,342\,015\,608$, le chiffre des millions est 2, mais le nombre de millions est $7\,342$ : on lit tout ce qui est écrit à gauche de ce rang, ce rang compris.\end{savoirfaire}
\essSection{2}{Unités, dixièmes, centièmes et millièmes}
\begin{propriete}Chaque rang vaut dix fois le rang situé à sa droite. Ainsi $1=\frac{10}{10}=\frac{100}{100}=\frac{1\,000}{1\,000}$ ; $\frac1{10}=\frac{10}{100}$ ; $\frac1{100}=\frac{10}{1\,000}$. Dans l’autre sens : $1=10\times\frac1{10}=100\times\frac1{100}$ et $\frac1{10}=10\times\frac1{100}$.\end{propriete}
\label{fig-cours-1-1-cours-centiemes}\figureGeo{cours-centiemes}

\begin{exemple}$\frac1{10}=0{,}1$ ; $\frac1{100}=0{,}01$ ; $\frac1{1\,000}=0{,}001$. Trois dixièmes et quatre centièmes s’écrivent $0{,}34$ : c’est aussi $\frac{34}{100}$, trente-quatre centièmes.\end{exemple}
\clearpage
\phantomsection\label{cours-1-2}
\essSection{3}{Reconnaître un nombre décimal}
\begin{definition}Un nombre décimal est un nombre qui peut s’écrire avec une fraction décimale, c’est-à-dire une fraction de dénominateur 10, 100, $1\,000$… Par exemple $2{,}5=\frac{25}{10}$ ; $\frac34=\frac{75}{100}$ ; $7=\frac{70}{10}$. Tout nombre entier est donc aussi un nombre décimal.\end{definition}
\begin{attention}Un nombre décimal n’est pas fait de deux entiers séparés par une virgule : dans $3{,}12$, la partie 12 compte des centièmes. Certaines fractions ne sont pas des nombres décimaux, comme $\frac13$ : le chapitre 5 explique pourquoi.\end{attention}
\essSection{4}{Passer d’une écriture à une autre}
\begin{propriete}Un même nombre a plusieurs écritures : $4{,}107=\frac{4\,107}{1\,000}=4+\frac{107}{1\,000}=4+\frac1{10}+\frac7{1\,000}$. L’écriture $4+\frac{107}{1\,000}$, un entier plus une fraction plus petite que 1, s’appelle un nombre mixte.\end{propriete}
\begin{savoirfaire}[Lire le dernier rang]Pour passer de la virgule à la fraction, on regarde le rang du dernier chiffre : $0{,}045$ se termine aux millièmes, donc $0{,}045=\frac{45}{1\,000}$. Dans l’autre sens, $\frac{2\,375}{100}$ compte des centièmes : $\frac{2\,375}{100}=23{,}75$.\end{savoirfaire}
\begin{exemple}Un zéro écrit à droite de la partie décimale ne change pas le nombre : $2{,}5=2{,}50=2{,}500$, car 5 dixièmes font 50 centièmes, ou 500 millièmes.\end{exemple}

\clearpage\rubrique{Découverte}
\begin{blocLivre}{ex-1-1}\begin{exercice}[1]{Lire un grand nombre}
Écrire en lettres le nombre $3\,052\,400\,000$, puis écrire en chiffres « douze-milliards-cinq-cents-millions ».
\end{exercice}\end{blocLivre}
\begin{blocLivre}{ex-1-2}\begin{exercice}[1]{Chiffre ou nombre}
Dans $7\,342\,015\,608$, donner le chiffre des millions et le nombre de millions, puis le chiffre des centaines de mille et le nombre de centaines de mille.
\end{exercice}\end{blocLivre}
\begin{blocLivre}{ex-1-3}\begin{exercice}[1]{Unités de numération}
Compléter : $1=\square$ dixièmes ; 1 dixième $=\square$ centièmes ; 1 centième $=\square$ millièmes ; $1=\square$ millièmes ; 1 millier $=\square$ centaines.
\end{exercice}\end{blocLivre}
\begin{blocLivre}{ex-1-4}\begin{exercice}[1]{Lire une grille}
L’unité est un grand carré partagé en cent petits carrés. Écrire le nombre représenté par la partie colorée sous forme de fraction décimale, puis avec une virgule.
\label{fig-ex-1-4}\figureGeo{num-4}
\end{exercice}\end{blocLivre}
\begin{blocLivre}{ex-1-5}\begin{exercice}[1]{Fractions décimales}
Écrire avec une virgule : $\frac7{10}$ ; $\frac{45}{100}$ ; $\frac3{1\,000}$ ; $4+\frac1{10}+\frac7{1\,000}$.
\end{exercice}\end{blocLivre}

\rubrique{Entraînement}
\begin{blocLivre}{ex-1-6}\begin{exercice}[2]{Valeur des chiffres}
Dans le nombre $58{,}036$, donner la valeur de chacun des chiffres.
\end{exercice}\end{blocLivre}
\begin{blocLivre}{ex-1-7}\begin{exercice}[2]{Le même nombre ?}
Parmi ces écritures, lesquelles désignent le nombre $3{,}4$ ? $\frac{34}{10}$ ; $3+\frac4{100}$ ; $\frac{340}{100}$ ; $3+\frac4{10}$ ; $\frac{34}{100}$.
\end{exercice}\end{blocLivre}
\begin{blocLivre}{ex-1-8}\begin{exercice}[2]{Nombres mixtes}
Écrire chaque nombre comme un entier plus une fraction décimale plus petite que 1, puis avec une virgule : $\frac{2\,375}{1\,000}$ ; $\frac{1\,204}{100}$ ; $\frac{89}{10}$.
\end{exercice}\end{blocLivre}
\begin{blocLivre}{ex-1-9}\begin{exercice}[2]{Décimal ou pas ?}
Montrer que chacun de ces nombres est décimal en l’écrivant avec une fraction décimale, puis avec une virgule : $\frac34$ ; $7$ ; $\frac15$ ; $\frac2{25}$.
\end{exercice}\end{blocLivre}
\begin{blocLivre}{ex-1-10}\begin{exercice}[2]{Distances dans l’espace}
La distance moyenne de la Terre au Soleil est d’environ 150 millions de kilomètres ; celle de Neptune au Soleil, d’environ $4{,}5$ milliards de kilomètres. Écrire ces deux distances en chiffres, puis trouver combien de fois la première est contenue dans la seconde.
\end{exercice}\end{blocLivre}

\rubrique{Problèmes}
\begin{blocLivre}{ex-1-11}\begin{exercice}[3]{Huit milliards d’humains}
Selon l’ONU, la population mondiale a atteint 8 milliards de personnes en novembre 2022 ; elle était de 4 milliards en 1974. Écrire ces deux nombres en chiffres. Combien de millions de personnes de plus en 2022 ? Combien de fois plus ?
\end{exercice}\end{blocLivre}
\begin{blocLivre}{ex-1-12}\begin{exercice}[3]{Des boîtes de mille}
Une usine fabrique $2\,450\,000\,000$ vis par an. Elle les range par boîtes de $1\,000$ vis, puis place les boîtes dans des cartons de $1\,000$ boîtes. Combien de boîtes remplit-elle ? Combien de cartons ?
\end{exercice}\end{blocLivre}
\begin{blocLivre}{ex-1-13}\begin{exercice}[3]{Le nombre mystère}
Je suis un nombre décimal. J’ai 3 unités et aucun centième ; mon chiffre des dixièmes est 4 et mon chiffre des millièmes est le double de celui des dixièmes. Je n’ai pas d’autre chiffre non nul. Qui suis-je ? Donner aussi mon écriture en fraction décimale.
\end{exercice}\end{blocLivre}
\begin{blocLivre}{ex-1-14}\begin{exercice}[3]{La notation de Stevin}
En 1585, le savant Simon Stevin écrivait 27\stevin{0}8\stevin{1}4\stevin{2}7\stevin{3} pour $27{,}847$ : le chiffre entouré indique le rang (0 pour les unités, 1 pour les dixièmes, 2 pour les centièmes…). Écrire $5{,}306$ avec sa notation, puis écrire avec une virgule le nombre 3\stevin{0}0\stevin{1}9\stevin{2}.
\end{exercice}\end{blocLivre}
\begin{blocLivre}{ex-1-15}\begin{exercice}[3]{Vrai ou faux}
Léa affirme : « Dans $6{,}25$, il y a 25 dixièmes. » A-t-elle raison ? Combien de dixièmes, puis combien de centièmes, le nombre $6{,}25$ contient-il ?
\end{exercice}\end{blocLivre}

\clearpage
